\documentclass[aspectratio=169]{beamer}

\usetheme{Madrid}
\usecolortheme{default}
\usepackage{booktabs}
\usepackage{graphicx}
\usepackage{amsmath}
\usepackage[T1]{fontenc}
\usepackage[utf8]{inputenc}

% Custom colors
\definecolor{mainblue}{RGB}{0,70,140}
\setbeamercolor{palette primary}{bg=mainblue,fg=white}
\setbeamercolor{palette secondary}{bg=mainblue!80,fg=white}
\setbeamercolor{palette tertiary}{bg=mainblue!60,fg=white}
\setbeamercolor{palette quaternary}{bg=mainblue,fg=white}
\setbeamercolor{structure}{fg=mainblue}
\setbeamercolor{section in toc}{fg=mainblue}

\setbeamertemplate{navigation symbols}{}
\setbeamertemplate{footline}[frame number]

% -------------------------------------------------------
\title{SSL Pretraining Study with Primate Data\\for Human Speech Decoding}
\author{Sergio López}
\date{}
% -------------------------------------------------------

\begin{document}

\begin{frame}
  \titlepage
\end{frame}

% ---- OUTLINE -------------------------------------------
\begin{frame}{Outline}
  \tableofcontents
\end{frame}

% ========================================================
\section{Introduction}
% ========================================================

\begin{frame}{Introduction — Kaggle Brain-to-Text '25}
  \begin{columns}
    \begin{column}{0.55\textwidth}
      \begin{itemize}
        \item \textbf{Task:} Speech BCI to restore communication of paralysed patients
        \item \textbf{Goal:} Neural Activity $\rightarrow$ Text
        \item \textbf{Competition metric:} Word Error Rate (WER)
        \item \textbf{Our focus:} Phoneme Error Rate (PER)
      \end{itemize}
    \end{column}
    \begin{column}{0.4\textwidth}
      % Insert competition overview figure here
      \begin{center}
        \fbox{\parbox{0.9\textwidth}{\centering\small [Figure: BCI overview]}}
      \end{center}
    \end{column}
  \end{columns}
\end{frame}

% ========================================================
\section{Dataset}
% ========================================================

\begin{frame}{Dataset — Participant T15}
  \begin{columns}
    \begin{column}{0.55\textwidth}
      \begin{itemize}
        \item $\sim$11\,000 sentences recorded over 20 months
        \item 256 microelectrodes in speech motor cortex
        \item Time bins: 20\,ms (50\,Hz)
        \item Multiple sentence corpora:
          \begin{itemize}
            \item 50-word vocabulary
            \item Switchboard, OpenWebText
            \item Harvard sentences, random words
          \end{itemize}
        \item Two speaking strategies:
          \begin{itemize}
            \item Vocalized ($\sim$30 wpm)
            \item Silent ($\sim$50 wpm)
          \end{itemize}
      \end{itemize}
    \end{column}
    \begin{column}{0.4\textwidth}
      \begin{center}
        \fbox{\parbox{0.9\textwidth}{\centering\small [Figure: electrode array / recording setup]}}
      \end{center}
    \end{column}
  \end{columns}
\end{frame}

% ========================================================
\section{Competition Work}
% ========================================================

\begin{frame}{My Competition Work — Improvements to the Baseline}
  \begin{columns}
    \begin{column}{0.5\textwidth}
      \textbf{Optimisations applied:}
      \begin{itemize}
        \item Cosine LR scheduling
        \item Speckled masking
        \item Hyperparameter tuning:
          \begin{itemize}
            \item Dropout: 0.2
            \item Learning rate: 0.004
            \item Batch size: 32
          \end{itemize}
      \end{itemize}
      \vspace{0.5em}
      \begin{block}{Result}
        \centering $\downarrow$ PER = \textbf{0.097}
      \end{block}
    \end{column}
    \begin{column}{0.45\textwidth}
      \begin{center}
        \fbox{\parbox{0.9\textwidth}{\centering\small [Figure: training curves]}}
      \end{center}
    \end{column}
  \end{columns}
\end{frame}

% ========================================================
\section{Motivation: The BIT Paper}
% ========================================================

\begin{frame}{Motivation — The BIT Paper}
  \begin{columns}
    \begin{column}{0.55\textwidth}
      \textbf{BIT (Ye et al., 2025):}
      \begin{itemize}
        \item Pretrained transformer with human + NHP data
        \item PER of \textbf{7.12\%} on T15
      \end{itemize}
      \vspace{0.8em}
      \textbf{Non-human primate (NHP) data:}
      \begin{itemize}
        \item Abundant and publicly available
        \item Recorded with Utah array technology
      \end{itemize}
      \vspace{0.8em}
      \begin{alertblock}{Research Question}
        Can we improve a human speech decoder using \textbf{only NHP data} for pretraining?
      \end{alertblock}
    \end{column}
    \begin{column}{0.4\textwidth}
      \begin{center}
        \fbox{\parbox{0.9\textwidth}{\centering\small [Figure: BIT architecture / NHP illustration]}}
      \end{center}
    \end{column}
  \end{columns}
\end{frame}

% ========================================================
\section{NHP Datasets}
% ========================================================

\begin{frame}{Non-Human Primate Datasets}
  \begin{columns}
    \begin{column}{0.55\textwidth}
      \textbf{200 hours of NHP data} across tasks:
      \begin{itemize}
        \item Reaching
        \item Fine motor
        \item Somatosensory
      \end{itemize}
      \vspace{0.5em}
      \textbf{Preprocessing (from BIT):}
      \begin{itemize}
        \item 20\,ms time bins
        \item Z-scoring across days
      \end{itemize}
    \end{column}
    \begin{column}{0.4\textwidth}
      \begin{center}
        \fbox{\parbox{0.9\textwidth}{\centering\small [Figure: NHP dataset overview]}}
      \end{center}
    \end{column}
  \end{columns}
\end{frame}

% ========================================================
\section{GRU SSL Pretraining}
% ========================================================

\begin{frame}{GRU SSL Pretraining — Method}
  \begin{columns}
    \begin{column}{0.5\textwidth}
      \textbf{Pretraining (on NHP data):}
      \begin{itemize}
        \item Causal next-step prediction
        \item MSE loss on next 3 patches
      \end{itemize}
      \vspace{0.5em}
      \textbf{Transfer to human data:}
      \begin{itemize}
        \item GRU weights transferred
        \item Day-specific input layer reinitialized
        \item CTC head reinitialized
      \end{itemize}
    \end{column}
    \begin{column}{0.45\textwidth}
      \begin{center}
        \fbox{\parbox{0.9\textwidth}{\centering\small [Figure: GRU pretraining diagram]}}
      \end{center}
    \end{column}
  \end{columns}
\end{frame}

\begin{frame}{GRU SSL Pretraining — Results}
  \begin{columns}
    \begin{column}{0.5\textwidth}
      \begin{table}
        \centering
        \begin{tabular}{lc}
          \toprule
          \textbf{Dataset} & \textbf{PER} \\
          \midrule
          GRU control (no SSL) & 0.097 \\
          Somatosensory        & 0.724 \\
          Fine motor           & 0.718 \\
          Reaching             & 0.140 \\
          All NHP              & 0.136 \\
          \bottomrule
        \end{tabular}
      \end{table}
    \end{column}
    \begin{column}{0.45\textwidth}
      \begin{alertblock}{Why SSL fails on GRU}
        Recurrent dynamics are coupled to input distribution — transfer breaks when going NHP $\to$ human.
      \end{alertblock}
      \vspace{0.4em}
      \small
      Churchland et al.\ 2012 (``Neural population dynamics during reaching'', \textit{Nature}) \\
      BIT (Ye et al.\ 2025)
    \end{column}
  \end{columns}
\end{frame}

% ========================================================
\section{Transformer Approach}
% ========================================================

\begin{frame}{Building Our Transformer}
  \begin{columns}
    \begin{column}{0.5\textwidth}
      \textbf{BIT transformer:}
      \begin{itemize}
        \item Works with SSL
        \item Beats GRU baseline
        \item NHP $\to$ human transfer works
      \end{itemize}
      \vspace{0.5em}
      \textbf{Our additions:}
      \begin{itemize}
        \item Day-specific input layer
        \item Aggressive augmentation (noise=0.8, offset=0.2)
        \item Log-transform on input features
        \item HP tuning: lr=$3\times10^{-4}$, WD=$5\times10^{-4}$, batch=32
      \end{itemize}
      \vspace{0.4em}
      \begin{block}{Result}
        \centering $\downarrow$ PER = \textbf{0.097} (matches optimised GRU)
      \end{block}
    \end{column}
    \begin{column}{0.45\textwidth}
      \begin{center}
        \fbox{\parbox{0.9\textwidth}{\centering\small [Figure: transformer architecture]}}
      \end{center}
    \end{column}
  \end{columns}
\end{frame}

\begin{frame}{SSL Pretraining with Transformer}
  \begin{columns}
    \begin{column}{0.5\textwidth}
      \begin{itemize}
        \item Pretraining loss and $R^2$ do \textbf{not} predict PER
        \item Evaluated every few pretraining epochs with short fine-tuning on human data
      \end{itemize}
      \vspace{0.5em}
      \textbf{PER at intermediate checkpoints:}
      \begin{itemize}
        \item Ep.\ 50:  0.169
        \item Ep.\ 100: 0.160
        \item Ep.\ 150: 0.146
        \item Ep.\ 200: 0.152
        \item Ep.\ 250: 0.165
      \end{itemize}
    \end{column}
    \begin{column}{0.45\textwidth}
      \begin{center}
        \fbox{\parbox{0.9\textwidth}{\centering\small [Figure: pretraining loss vs.\ PER curves]}}
      \end{center}
    \end{column}
  \end{columns}
\end{frame}

\begin{frame}{SSL Study — PER at 30k Batches}
  \begin{table}
    \centering
    \begin{tabular}{lcccc}
      \toprule
      \textbf{Dataset} & \textbf{Hours} & \textbf{Masked SSL} & \textbf{AR binary} \\
      \midrule
      Control (no SSL)  & ---   & 0.124              & 0.124 \\
      Reaching          & 117h  & 0.142              & 0.129 \\
      Fine motor        & 68h   & 0.203              & 0.437 \\
      Somatosensory     & 10h   & 0.124 (= control)  & \textbf{0.118} \\
      All NHP           & 200h  & 0.137              & 0.138 \\
      \bottomrule
    \end{tabular}
    \caption{PER values after fine-tuning on T15 (30k batches)}
  \end{table}
\end{frame}

% ========================================================
\section{Best Result}
% ========================================================

\begin{frame}{Best Result}
  \begin{columns}
    \begin{column}{0.55\textwidth}
      \begin{itemize}
        \item \textbf{Best checkpoint:} AR binary somatosensory, epoch 400
        \item Full fine-tuning: 400k batches on T15
      \end{itemize}
      \vspace{0.6em}
      \begin{block}{Final PER}
        \centering \Large PER = \textbf{0.0887}
      \end{block}
      \vspace{0.6em}
      \begin{exampleblock}{Improvement}
        \centering \textbf{8.6\% improvement} over TFS baseline (0.097)
      \end{exampleblock}
    \end{column}
    \begin{column}{0.4\textwidth}
      \begin{center}
        \fbox{\parbox{0.9\textwidth}{\centering\small [Figure: final fine-tuning curve]}}
      \end{center}
    \end{column}
  \end{columns}
\end{frame}

% ========================================================
\section{Next Steps}
% ========================================================

\begin{frame}{Possible Next Steps}
  \begin{columns}
    \begin{column}{0.48\textwidth}
      \textbf{Verify robustness:}
      \begin{itemize}
        \item Is the gain due to dataset size?
        \item Multi-seed fine-tuning to confirm best result is not seed-dependent
      \end{itemize}
      \vspace{0.5em}
      \textbf{Expand SSL comparison:}
      \begin{itemize}
        \item Additional SSL techniques on the full dataset grid
        \item Combination of NHP groups
        \item Other cortical areas?
      \end{itemize}
    \end{column}
    \begin{column}{0.48\textwidth}
      \textbf{Deeper analysis:}
      \begin{itemize}
        \item Why does \textit{somatosensory} cortex transfer but motor cortex does not?
      \end{itemize}
    \end{column}
  \end{columns}
\end{frame}

% ========================================================
\begin{frame}
  \centering
  \Large\textbf{Thank you!}\\[1em]
  \normalsize Questions?
\end{frame}

\end{document}
